\section{The native potential of the same original tensor power}

\begin{theorem}[Actual native norm factor of the original power data]
\label{thm:native-tc-nef-fs-power-factor}
\leanok
\uses{lem:native-tc-nef-fs-native-factor-arbitrary-center, lem:native-tc-nef-fs-product-power-coherence}
\lean{EnvelopeProperty.NativeTestConfiguration.TestConfiguration.gaussNativeNormFactor_power_comparisons}
Fix the same original test configuration, line $M$, generic line $L$,
whole action comparison $\alpha$ and fibre-one comparison $e_1$.
If $L^{\otimes n}$ is actually ample on the original $X$, then the native
norm factor for the internally constructed original power data equals
the $n$th power of the original native norm factor, at every original
native point and radius.
\end{theorem}
\begin{proof}
Use the same original proper-centre open and frame for $L$ and their
actual $n$th power for $L^{\otimes n}$.  At the original total-model
valuation centre use the $n$th centre power frame of the original
chosen centre frame.  The arbitrary-centre-frame factor theorem
identifies the native factor of the powered data with this factor;
its ampleness hypothesis is exactly the stated actual ampleness of
$L^{\otimes n}$.  Whole power-action and fibre-one coherence identify
the induced original product model frame with the original canonical
product power frame.  The original power-frame transition theorem then
identifies the actual transition unit with the $n$th power of the
original transition unit.  The actual residue valuation preserves
powers.  The degree-one transition factor is literally the defined
original native factor, with its original chosen open frame.
\end{proof}

\begin{theorem}[Actual original native potential power law]
\label{thm:native-tc-nef-fs-power-potential}
\leanok
\uses{thm:native-tc-nef-fs-power-factor, def:native-tc-original-model-potential}
\lean{EnvelopeProperty.NativeTestConfiguration.TestConfiguration.gaussNativeModelPotential_power_comparisons}
Under the same hypotheses, the native model potential of the actual
power data is $n$ times the original native model potential.
\end{theorem}
\begin{proof}
The native potential is minus the real logarithm of the positive native
factor.  Apply the preceding exact factor power law and the real
logarithm power identity.
\end{proof}
